\section{CONCLUSÕES E TRABALHOS FUTUROS}
\label{sec:conclusao}

Neste trabalho, relatou-se a experiência de análise e reestruturação da documenta-\\ção de requisitos e de testes no projeto SFM. Ao longo da pesquisa, fundamentada pela construção de uma Matriz de Rastreabilidade, diagnosticou-se uma cobertura de testes deficitária e observaram-se os impactos negativos das DTs acumuladas. Evidenciou-se que a aplicação rigorosa de práticas de Engenharia de Software é essencial para resguardar a qualidade do produto. Notou-se, também, a necessidade de uma frente dedicada à Engenharia de Requisitos, dada a inviabilidade de sobrecarregar outras equipes — como a de \textit{design} — com essa atribuição. Sobretudo, constatou-se que a documentação formal possui importância vital não apenas para a manutenção do sistema por futuras equipes, mas para o seu contínuo desenvolvimento e aprimoramento.

Dessa forma, a presente pesquisa contribuiu de maneira prática por meio da entrega de uma documentação de requisitos que mapeia e abrange as reais necessidades do SFM. Além do artefato concebido, este estudo atua como um alerta para equipes de desenvolvimento corporativo em geral: a negligência documental nos estágios iniciais é altamente prejudicial, e os ``juros'' decorrentes do acúmulo de DTs inevitavelmente demandarão mitigação em algum momento do ciclo. De maneira mais específica para a realidade do projeto avaliado, o trabalho evidencia os desafios inerentes a um fluxo de desenvolvimento empírico, no qual a ausência de artefatos formais resultou em retrabalhos estruturais passíveis de mitigação e em uma má compreensão do estado do sistema implementado. 

Como desdobramento prático e operacional, recomenda-se fortemente que a coordenação do projeto incorpore os artefatos elaborados ao seu fluxo de trabalho. A documentação estabelecida deve passar a atuar como um critério restritivo, de modo a garantir que nenhuma nova funcionalidade seja codificada sem a sua devida especificação. Nesse sentido, embora a versão inicial do artefato já tenha sido homologada pela gestão, a natureza evolutiva do sistema exige que a equipe assuma a responsabilidade por sua manutenção e validação contínuas. À medida que novas demandas são debatidas nas reuniões de alinhamento, torna-se imperativo atualizar e revalidar o documento, além de inspecionar os módulos já implementados, com o intuito de assegurar a exatidão e a vigência das informações. É imprescindível que qualquer alteração no código seja imediatamente refletida nos requisitos, para que a especificação não se torne obsoleta e continue atuando como uma fonte de verdade confiável.

Adicionalmente, almeja-se que a documentação atue simultaneamente como base estrutural para a automação dos testes de \textit{software}. Como explanado previamente, a dependência exclusiva da execução manual de testes torna-se insustentável ao longo do ciclo de desenvolvimento, uma vez que restringe o potencial estratégico da equipe de QA. Portanto, recomenda-se que os comportamentos já mapeados nos casos de uso sejam transpostos para uma sintaxe de especificação executável, a exemplo do \textit{Gherkin}. Essa abordagem permite converter a linguagem natural da documentação em instruções interpretáveis por \textit{frameworks} de automação. Dessa forma, a validação do sistema receberá um suporte mais robusto do setor de qualidade, provendo casos de teste capazes de refletir fielmente a especificação, com o intuito de assegurar a confiabilidade e a integridade do produto.

Por fim, torna-se imperativa a apresentação das limitações inerentes a este estudo. Primariamente, pontua-se a natureza mutável do \textit{software}. Uma vez que todo ecossistema computacional é orgânico e iterativo, a análise aqui exposta reflete um recorte temporal específico, correspondente ao período em que a equipe enfrentava os gargalos estruturais relatados. Em segundo lugar, o escopo da intervenção restringiu-se à reestruturação da documentação funcional e de testes, e o instrumento de coleta de dados limitou-se a aferir a percepção técnica sobre o estado do sistema e a usabilidade do artefato produzido. Nesse sentido, variáveis externas não foram objeto de investigação, tais como o uso de documentações legadas, o impacto de diferentes formatos textuais na compreensão da equipe, e o nível de preferência por artefatos integrados a ambientes de desenvolvimento e versionamento (como o \textit{GitHub}). Adicionalmente, no que tange à aplicação desse formulário, pontua-se a ausência de definições operacionais estritas para os termos avaliados. Conceitos-chave, a exemplo da métrica de "\textit{software} pronto", bem como a semântica exata dos graus da escala Likert adotada, não foram previamente padronizados, sujeitando as respostas à interpretação subjetiva de cada respondente. Constata-se, ainda, a restrição da amostra à equipe de desenvolvimento, não englobando a perspectiva gerencial da coordenação. Por fim, a especificação absteve-se de mapear os requisitos não funcionais, os quais são cruciais para viabilizar outros vieses de teste no futuro.

No que tange às perspectivas para trabalhos acadêmicos futuros, sugere-se, primariamente, a expansão do escopo documental por meio do mapeamento e da formalização dos requisitos não funcionais do SFM. Propõe-se, também, a condução de um estudo longitudinal no projeto, com o intuito de elaborar uma análise comparativa capaz de mensurar os impactos reais na produtividade da equipe e na qualidade do \textit{software} após a adoção consolidada dessas práticas. Recomenda-se, adicionalmente, investigar a influência do formato e do local de hospedagem dos artefatos (sejam eles documentos na nuvem ou repositórios de código) na adesão técnica e no fluxo de trabalho diário. Sob a ótica administrativa, a realização de um levantamento focado nas perspectivas gerenciais elucidaria como as DTs são identificadas, priorizadas e geridas pelas instâncias superiores do projeto. Por fim, incentiva-se a condução de análises direcionadas à manutenibilidade do sistema, buscando dimensionar como a presença de uma especificação formal atua na mitigação de riscos associados à rotatividade de membros da equipe e à evolução contínua da plataforma.